% Experiment report. Fixed sections: see experiments/TEMPLATE/report.tex and experiments/README.md.
% Build: npm run experiment:report -- TEMPLATE
\documentclass{../polyexp}
\begin{document}
% generated by scripts/experiment-report.ts from experiment.json and models/<run>/ - do not edit
\providecommand{\ExpHeader}{}\renewcommand{\ExpHeader}{%
{\sffamily\footnotesize\color{expgrey}EXPERIMENT 17 \textperiodcentered\ TUNING \hfill 2026-11-26}\par
{\LARGE\bfseries\sffamily Fine-tune the gate's greedy bots for Perfection}\hfill\verdictbox{expgreen}{KEPT}\par\smallskip
{\itshape Does a perturbation-probe tuning run in Perfection (Continents, 30 turns, share of points) make the gate's two greedy champions stronger in Perfection than the Domination-tuned originals?}\par\medskip
\begin{tabular}{@{}p{22mm}p{\dimexpr\linewidth-26mm}@{}}
\fact{Change}{No new mechanism: the training mode changes from Domination/Drylands to Perfection/Continents, and the metric from the finish score to the share of points. Two runs, one per greedy model in the gate: probe (champion and frozen anchor) and combo (previous champion).}
\fact{Base}{\model{greedy-full-probe} $\rightarrow$ result \model{greedy-full-probe-perf passes rl-gate vs probe (+176 mean, 53.5\% of seeds better) and becomes champion and frozen anchor; greedy-full-combo-perf closes combo's -415 gap to +16 and becomes previous champion (Jakob's decision)}}
\fact{Method}{Perturbation probe, 100k games per run, 10k games per round, sigma 0.25, z 3, step 0.5, league = the gate's opponents (greedy-full-probe, greedy-full-combo); 400-game final test against each panel model; metric: share of points}
\fact{Setups}{Perfection, Crazy, 4 players, 16$\times$16, continents, 30 turns, no human seat (trainer); then the rl-gate setup (Crazy, learner in the human seat) as a check}
\fact{Record}{issue \issue{166}; status: done}
\end{tabular}\par\medskip
\begin{center}\begin{tikzpicture}[node distance=5mm]
    \node[flowbox] (n0) {{\sffamily\bfseries\scriptsize\color{expgrey}BASE}\\[1pt]\model{full-probe}};
    \node[flowbox,right=5mm of n0] (n1) {{\sffamily\bfseries\scriptsize\color{expgrey}CHANGE}\\[1pt]tuning};
    \node[flowbox,right=5mm of n1] (n2) {{\sffamily\bfseries\scriptsize\color{expgrey}MEASURED}\\[1pt]Panel curve every round and the final test, in Perfection\\{\scriptsize 1 setup, metric: share of points}};
    \node[flowbox,right=5mm of n2] (n3) {{\sffamily\bfseries\scriptsize\color{expgrey}RESULT}\\[1pt]Share of points vs the base, final test (probe run)\\{\scriptsize 50.0 $\rightarrow$ 53.6\,\%}\\\textbf{+3.6\pp}};
    \node[flowbox,right=5mm of n3] (n4) {{\sffamily\bfseries\scriptsize\color{expgrey}VERDICT}\\[1pt]\textbf{\color{expgreen}KEPT}\\{\scriptsize new champion:}\\\model{full-probe-perf passes rl-gate vs probe (+176 mean, 53.5\% of seeds better) and becomes champion and frozen anchor; greedy-full-combo-perf closes combo's -415 gap to +16 and becomes previous champion (Jakob's decision)}};
    \draw[flowarrow] (n0) -- (n1); \draw[flowarrow] (n1) -- (n2); \draw[flowarrow] (n2) -- (n3); \draw[flowarrow] (n3) -- (n4);
\end{tikzpicture}\end{center}
\begin{center}\sffamily\small\begin{tabular}{@{}lrrr@{}}\toprule
Measure & Base & New & Change\\\midrule
Share of points vs the base, final test (probe run) & 50.0\,\% & 53.6\,\% & +3.6\\
Share of points vs the base, final test (combo run) & 50.0\,\% & 56.7\,\% & +6.7\\
\bottomrule\end{tabular}\end{center}
}
\providecommand{\ExpResults}{}\renewcommand{\ExpResults}{%
\subsection{Probe run: greedy-full-probe tuned for Perfection {\normalfont\small(\model{greedy-full-probe-perf})}}
\begin{center}\begin{tikzpicture}\begin{axis}[expplot,height=4.6cm,xlabel={training games (thousands)},ylabel={Points share, panel mean (\%)},enlarge y limits=0.2]
\addplot[dashed,expgrey,line width=0.7pt,forget plot] coordinates {(0.0,71.54) (100.0,71.54)};
\addplot[expblue,mark=*,mark size=1.8pt] coordinates {(0.0,71.54) (20.0,71.90) (40.0,70.97) (60.0,70.88) (80.0,72.10) (100.0,72.42)};
\end{axis}\end{tikzpicture}\\[-2pt]{\footnotesize\color{expgrey}Learning curve of the tuning run (34 kept changes in 10 rounds). Dashed: the start.}\end{center}
\begin{center}\sffamily\small\begin{tabular}{@{}lrrrrr@{}}\toprule
Opponent & Won (\%) & Points (\%) & Finish (\%) & Conquest (\%) & Turns\\\midrule
\model{v1} & 100 $\pm$ 0 & 82.0 & 91.3 & 92 & 48\\
\model{hill-v1} & 98 $\pm$ 1 & 79.7 & 89.1 & 89 & 50\\
\model{hill-v2} & 98 $\pm$ 2 & 78.7 & 89.2 & 90 & 51\\
\model{hill-v3} & 97 $\pm$ 2 & 76.6 & 88.6 & 89 & 51\\
\model{hill-v4} & 96 $\pm$ 2 & 76.1 & 88.4 & 91 & 48\\
\model{full-v1} & 80 $\pm$ 4 & 68.2 & 74.7 & 73 & 55\\
\model{full-probe} & 56 $\pm$ 5 & 53.6 & 55.6 & 49 & 59\\
\model{full-combo} & 67 $\pm$ 5 & 59.7 & 63.0 & 58 & 58\\
\midrule
\textbf{Mean without \model{full-probe}} & 90.7 & 74.4 & 83.5 & 83.2 & 51.6\\
\bottomrule\end{tabular}\\[3pt]{\footnotesize\color{expgrey}Final test of the saved model: 400 unseen games against each panel model.}\end{center}
\begin{center}\begin{tikzpicture}\begin{axis}[expplot,height=4.8cm,xmin=-5,xmax=5,ymin=0,xlabel={z-score of a knob (up games minus down games, in standard errors)},ylabel={knobs}]
\addplot[ybar interval,fill=expblue!55,draw=white] coordinates {(-5.0,0) (-4.5,0) (-4.0,0) (-3.5,3) (-3.0,5) (-2.5,4) (-2.0,9) (-1.5,13) (-1.0,28) (-0.5,26) (0.0,48) (0.5,30) (1.0,22) (1.5,11) (2.0,4) (2.5,1) (3.0,1) (3.5,0) (4.0,0) (4.5,0) (5.0,0)};
\addplot[dashed,expgrey,line width=0.8pt,domain=-5:5,samples=80] {205*0.5*exp(-x^2/2)/sqrt(2*pi)};
\draw[expgrey,line width=0.4pt] (axis cs:-3,0) -- (axis cs:-3,\pgfkeysvalueof{/pgfplots/ymax}) (axis cs:3,0) -- (axis cs:3,\pgfkeysvalueof{/pgfplots/ymax});
\end{axis}\end{tikzpicture}\\[-2pt]{\footnotesize\color{expgrey}z-scores of all 205 knobs after round 10 (100k training games). Dashed: pure chance. Vertical lines: the threshold $|z|=3$; 4 beyond it. The outer bars hold everything past $\pm 5$.}\end{center}
\begin{center}\sffamily\small\begin{tabular}{@{}rlrrr@{}}\toprule
Round & Knob & From & To & z\\\midrule
2 & \texttt{tech.philosophy} & 3.821 & 4.33 & $+3.6$\\
2 & \texttt{centre} & 0.778 & 0.687 & $-3.7$\\
3 & \texttt{type.build} & 1 & 0.882 & $-3.4$\\
3 & \texttt{tech.climbing} & 1.284 & 1.133 & $-3.3$\\
3 & \texttt{unit.defender} & 1 & 0.882 & $-3.3$\\
3 & \texttt{front} & 1.133 & 1.284 & $+3.4$\\
4 & \texttt{type.research} & 0.535 & 0.472 & $-4.1$\\
4 & \texttt{tech.freeSpirit} & 0.882 & 0.778 & $-3.2$\\
4 & \texttt{tech.climbing} & 1.133 & 1.284 & $+5.8$\\
4 & \texttt{tech.hunting} & 1 & 0.882 & $-3.4$\\
4 & \texttt{build.fertilize} & 1 & 1.133 & $+3.0$\\
5 & \texttt{unit.siren} & 1.319 & 1.495 & $+3.2$\\
5 & \texttt{build.parkOfFortune} & 0.972 & 1.101 & $+3.3$\\
6 & \texttt{tech.mining} & 0.998 & 1.131 & $+3.3$\\
6 & \texttt{tech.meditation} & 1.848 & 2.094 & $+3.5$\\
6 & \texttt{plan} & 2 & 2.266 & $+3.1$\\
7 & \texttt{tech.philosophy} & 4.33 & 4.907 & $+4.0$\\
7 & \texttt{unit.mermaid} & 1 & 0.882 & $-3.2$\\
7 & \texttt{build.lumberHut} & 1 & 0.882 & $-3.2$\\
8 & \texttt{type.research} & 0.472 & 0.417 & $-3.2$\\
8 & \texttt{type.build} & 0.882 & 0.778 & $-3.8$\\
8 & \texttt{build.hunt} & 1 & 1.133 & $+3.1$\\
8 & \texttt{build.temple} & 1 & 1.133 & $+3.3$\\
8 & \texttt{moveNoise} & 2.266 & 2.568 & $+3.9$\\
9 & \texttt{greedy.kill} & 27.66 & 24.41 & $-3.1$\\
9 & \texttt{greedy.special} & 42.73 & 48.419 & $+3.3$\\
9 & \texttt{tech.construction} & 1.348 & 1.19 & $-3.2$\\
9 & \texttt{unit.siren} & 1.495 & 1.694 & $+3.3$\\
9 & \texttt{build.port} & 0.695 & 0.613 & $-3.2$\\
9 & \texttt{centre} & 0.687 & 0.606 & $-3.6$\\
\bottomrule\end{tabular}\\[3pt]{\footnotesize\color{expgrey}Knobs moved: 34 (the first 30 shown; all in \texttt{probe.jsonl}).}\end{center}
\begin{center}\sffamily\small\begin{tabular}{@{}lrrrrrr@{}}\toprule
Knob & Value & Samples & Effect (pts) & z & z self & z league\\\midrule
\texttt{centre} & 0.606 & 1667 & $-2.44$ & $-3.3$ & $-3.1$ & $-2.2$\\
\texttt{tech.philosophy} & 4.907 & 5001 & $+1.39$ & $+3.3$ & $+1.7$ & $+3.2$\\
\texttt{tech.spiritualism} & 1 & 16670 & $-0.75$ & $-3.2$ & $-2.6$ & $-2.5$\\
\texttt{greedy.growPerPop} & 29.26 & 16670 & $-0.72$ & $-3.1$ & $-3.1$ & $-1.9$\\
\texttt{tech.diplomacy} & 1.264 & 16670 & $-0.69$ & $-2.9$ & $-2.8$ & $-1.9$\\
\texttt{unit.mermaidArcher} & 1 & 16670 & $+0.67$ & $+2.8$ & $+0.8$ & $+3.3$\\
\texttt{greedy.growBase} & 16.947 & 16670 & $-0.65$ & $-2.8$ & $-2.1$ & $-2.2$\\
\texttt{type.build} & 0.778 & 3334 & $-1.42$ & $-2.7$ & $-2.2$ & $-2.1$\\
\texttt{unit.livingIsland} & 1 & 16670 & $-0.60$ & $-2.6$ & $-1.0$ & $-2.8$\\
\texttt{type.research} & 0.417 & 3334 & $-1.32$ & $-2.5$ & $-2.6$ & $-1.5$\\
\bottomrule\end{tabular}\\[3pt]{\footnotesize\color{expgrey}Strongest evidence at the end of the run. Effect: outcome gained per sample when the knob is nudged up instead of down, halved; the outcome adds the self-play and the league comparison.}\end{center}

\subsection{Combo run: greedy-full-combo tuned for Perfection {\normalfont\small(\model{greedy-full-combo-perf})}}
\begin{center}\begin{tikzpicture}\begin{axis}[expplot,height=4.6cm,xlabel={training games (thousands)},ylabel={Points share, panel mean (\%)},enlarge y limits=0.2]
\addplot[dashed,expgrey,line width=0.7pt,forget plot] coordinates {(0.0,68.09) (100.0,68.09)};
\addplot[expblue,mark=*,mark size=1.8pt] coordinates {(0.0,68.09) (20.0,70.28) (40.0,71.70) (60.0,71.61) (80.0,71.73) (100.0,71.85)};
\end{axis}\end{tikzpicture}\\[-2pt]{\footnotesize\color{expgrey}Learning curve of the tuning run (65 kept changes in 10 rounds). Dashed: the start.}\end{center}
\begin{center}\sffamily\small\begin{tabular}{@{}lrrrrr@{}}\toprule
Opponent & Won (\%) & Points (\%) & Finish (\%) & Conquest (\%) & Turns\\\midrule
\model{v1} & 99 $\pm$ 1 & 81.2 & 90.1 & 90 & 49\\
\model{hill-v1} & 99 $\pm$ 1 & 80.2 & 90.3 & 91 & 50\\
\model{hill-v2} & 99 $\pm$ 1 & 79.3 & 90.4 & 92 & 51\\
\model{hill-v3} & 95 $\pm$ 2 & 75.5 & 87.6 & 88 & 52\\
\model{hill-v4} & 96 $\pm$ 2 & 75.9 & 88.0 & 90 & 52\\
\model{full-v1} & 82 $\pm$ 4 & 68.7 & 75.8 & 75 & 53\\
\model{full-probe} & 52 $\pm$ 5 & 51.3 & 52.0 & 46 & 59\\
\model{full-combo} & 63 $\pm$ 5 & 56.7 & 59.1 & 54 & 58\\
\midrule
\textbf{Mean without \model{full-probe}} & 90.3 & 73.9 & 83.0 & 82.7 & 52.0\\
\bottomrule\end{tabular}\\[3pt]{\footnotesize\color{expgrey}Final test of the saved model: 400 unseen games against each panel model.}\end{center}
\begin{center}\begin{tikzpicture}\begin{axis}[expplot,height=4.8cm,xmin=-5,xmax=5,ymin=0,xlabel={z-score of a knob (up games minus down games, in standard errors)},ylabel={knobs}]
\addplot[ybar interval,fill=expblue!55,draw=white] coordinates {(-5.0,0) (-4.5,0) (-4.0,0) (-3.5,0) (-3.0,2) (-2.5,7) (-2.0,9) (-1.5,21) (-1.0,31) (-0.5,37) (0.0,36) (0.5,22) (1.0,18) (1.5,11) (2.0,5) (2.5,3) (3.0,3) (3.5,0) (4.0,0) (4.5,0) (5.0,0)};
\addplot[dashed,expgrey,line width=0.8pt,domain=-5:5,samples=80] {205*0.5*exp(-x^2/2)/sqrt(2*pi)};
\draw[expgrey,line width=0.4pt] (axis cs:-3,0) -- (axis cs:-3,\pgfkeysvalueof{/pgfplots/ymax}) (axis cs:3,0) -- (axis cs:3,\pgfkeysvalueof{/pgfplots/ymax});
\end{axis}\end{tikzpicture}\\[-2pt]{\footnotesize\color{expgrey}z-scores of all 205 knobs after round 10 (100k training games). Dashed: pure chance. Vertical lines: the threshold $|z|=3$; 3 beyond it. The outer bars hold everything past $\pm 5$.}\end{center}
\begin{center}\sffamily\small\begin{tabular}{@{}rlrrr@{}}\toprule
Round & Knob & From & To & z\\\midrule
1 & \texttt{greedy.research} & 27.6 & 31.275 & $+7.6$\\
1 & \texttt{type.research} & 1 & 0.882 & $-3.1$\\
1 & \texttt{tech.riding} & 0.948 & 0.837 & $-6.5$\\
1 & \texttt{tech.freeSpirit} & 1 & 0.882 & $-3.4$\\
1 & \texttt{tech.organization} & 1 & 1.133 & $+3.4$\\
1 & \texttt{tech.strategy} & 0.942 & 0.831 & $-4.0$\\
1 & \texttt{tech.climbing} & 1 & 1.133 & $+6.0$\\
1 & \texttt{tech.meditation} & 1.27 & 1.439 & $+5.9$\\
1 & \texttt{tech.philosophy} & 2.045 & 2.317 & $+7.1$\\
1 & \texttt{unit.warrior} & 1 & 1.133 & $+5.5$\\
1 & \texttt{unit.rider} & 1 & 0.882 & $-5.3$\\
1 & \texttt{reward.park} & 1 & 0.882 & $-5.0$\\
1 & \texttt{reward.superUnit} & 1 & 1.133 & $+7.0$\\
1 & \texttt{centre} & 1 & 0.882 & $-3.5$\\
2 & \texttt{type.research} & 0.882 & 0.778 & $-3.3$\\
2 & \texttt{tech.riding} & 0.837 & 0.739 & $-6.4$\\
2 & \texttt{tech.freeSpirit} & 0.882 & 0.778 & $-3.5$\\
2 & \texttt{tech.climbing} & 1.133 & 1.284 & $+3.7$\\
2 & \texttt{tech.meditation} & 1.439 & 1.631 & $+5.9$\\
2 & \texttt{tech.philosophy} & 2.317 & 2.626 & $+5.3$\\
2 & \texttt{tech.archery} & 1 & 0.882 & $-3.6$\\
2 & \texttt{tech.fishing} & 1 & 0.882 & $-4.1$\\
2 & \texttt{tech.sailing} & 0.923 & 0.815 & $-3.4$\\
2 & \texttt{unit.warrior} & 1.133 & 1.284 & $+7.0$\\
2 & \texttt{unit.rider} & 0.882 & 0.778 & $-5.1$\\
2 & \texttt{centre} & 0.882 & 0.778 & $-3.8$\\
3 & \texttt{greedy.special} & 42.73 & 48.419 & $+3.1$\\
3 & \texttt{type.research} & 0.778 & 0.687 & $-4.2$\\
3 & \texttt{tech.organization} & 1.133 & 1.284 & $+3.7$\\
3 & \texttt{tech.philosophy} & 2.626 & 2.976 & $+3.2$\\
\bottomrule\end{tabular}\\[3pt]{\footnotesize\color{expgrey}Knobs moved: 65 (the first 30 shown; all in \texttt{probe.jsonl}).}\end{center}
\begin{center}\sffamily\small\begin{tabular}{@{}lrrrrrr@{}}\toprule
Knob & Value & Samples & Effect (pts) & z & z self & z league\\\midrule
\texttt{unit.siren} & 1.319 & 3334 & $+1.76$ & $+3.4$ & $+3.2$ & $+2.3$\\
\texttt{greedy.special} & 54.866 & 5001 & $+1.37$ & $+3.2$ & $+3.0$ & $+2.1$\\
\texttt{tech.philosophy} & 4.33 & 3334 & $+1.64$ & $+3.1$ & $+3.3$ & $+1.8$\\
\texttt{type.attack} & 1.133 & 10002 & $+0.92$ & $+3.0$ & $+3.4$ & $+1.5$\\
\texttt{tech.spiritualism} & 1 & 16670 & $-0.71$ & $-2.8$ & $-2.2$ & $-2.3$\\
\texttt{type.research} & 0.417 & 3334 & $-1.43$ & $-2.7$ & $-3.0$ & $-1.6$\\
\texttt{unit.catapult} & 1 & 16670 & $+0.67$ & $+2.7$ & $+2.0$ & $+2.3$\\
\texttt{greedy.research} & 35.439 & 10002 & $+0.78$ & $+2.5$ & $+2.4$ & $+1.7$\\
\texttt{unit.exida} & 1 & 16670 & $+0.61$ & $+2.4$ & $+1.7$ & $+2.1$\\
\texttt{front} & 1.284 & 6668 & $+0.90$ & $+2.4$ & $+2.2$ & $+1.7$\\
\bottomrule\end{tabular}\\[3pt]{\footnotesize\color{expgrey}Strongest evidence at the end of the run. Effect: outcome gained per sample when the knob is nudged up instead of down, halved; the outcome adds the self-play and the league comparison.}\end{center}
}

\ExpHeader

\section{Question and hypothesis}
% WHEN: before the first run. WHO: Jakob states the idea in his words; Claude turns it into a prediction.
% Keep the prediction falsifiable: which number should move, in which direction, by about how much.
Jakob, 26 Nov 2026: ``I think we should use the current bots as they are as base version and then
just fine tune on top for perfection. They are already good at domination and there is some carry
over to perfection but perfection will need some adjustments due to the turn limit and objective
difference so finetuning is the way to go. However, we should only fine tune the ones in our gate
for now.'' The proposal came from David (\#166, 26 Nov): every greedy model was tuned only in
Domination on Drylands, and probe tuning even cut \texttt{type.research} to 0.535, while in
Perfection every tech adds score.

Prediction (written before the runs): each tuned model beats its base head to head in Perfection
with a share of points of at least 52\% in the 400-game final test, and \texttt{type.research}
rises above its base value (probe: 0.535). A share below 50.5\% for both runs, or a research
weight that falls, counts against the idea that Domination tuning is what holds the bots back in
Perfection.

\section{Mechanism}
No new mechanism and no code change. Only the training setup changes: \texttt{--mode perfection
--map-type continents --difficulty crazy --metric share} instead of Domination, Drylands, Normal
and the finish score. All knobs are free to move under the perturbation probe (experiment 12).
The trainer had no difficulty option: no greedy model had ever played Crazy (+4 capital income
for bots). The new \texttt{--difficulty} flag (\texttt{scripts/lib/train-greedy.ts}) passes it to
\texttt{createGame}; every trainer seat is a bot, so all seats get the bonus, as the gate's
opponent seats do. Check: 20 seeds, probe against itself, mean total score 32,979 on Normal and
41,684 on Crazy.

\section{Setup}
Two runs from the gate's greedy models (\texttt{docs/rl-project/champions.json}): the champion
\texttt{greedy-full-probe} (also the frozen anchor, which itself stays unchanged) and the previous
champion \texttt{greedy-full-combo}. Each: 100k games, rounds of 10k, $\sigma = 0.25$, $z = 3$,
step 0.5, 5 workers. The league is the gate's opponent set (probe, combo) instead of the
Domination league, so the bot is tuned against the opponents it meets in the gate. Panel: the
fixed six plus probe and combo, all in Perfection; means are compared on the fixed six. The
trainer plays two bot sides on alternating seats with no human seat, so a tuned model is checked
afterwards on the rl-gate setup (Crazy, learner in the human seat). The base row at round 0 of each
curve is the control: the change is the training mode, so a run without it would only repeat
experiment 12.

\section{Results}
% WHEN: after the runs. The tables and plots below are generated from the measured data; add two to four sentences that read them.
\ExpResults

\todo{claude}{Read the tables: the best point, the size of the effect against the noise (about $\pm$1.2 points on a 6-model mean), the head-to-head.}

\section{Behaviour}
% WHEN: after the runs. WHO: both. One count the mechanism should change, plus what was seen in watched games.
\todo{claude}{The behaviour metric: one count the mechanism should change, with base and new value.}
\todo{jakob}{What you saw in the live games, if you watched any.}

\section{Interpretation}
% WHEN: after the runs. WHO: both. Why it worked or did not; was the prediction right?
\todo{claude}{Why the result came out this way, and whether the prediction held.}
\todo{jakob}{Your reading: agree, disagree, what is missing?}

\section{Threats to validity}
% WHO: Claude. What could make this result wrong or narrow: panel, map sizes, sample size, interactions with other knobs.
\todo{claude}{Two or three concrete threats.}

\section{Verdict and next step}
% WHO: Claude proposes by the keep rules (.claude/skills/model-experiment/SKILL.md); Jakob decides anything outside them.
\todo{claude}{Kept or not, under which keep rule, the champion afterwards, and the next experiment this suggests.}

\section{Notes for the paper}
% WHO: Jakob, with Claude's draft. What this experiment contributes: a claim, a figure, a negative result worth a line.
\todo{claude}{Draft: the one-sentence claim and which figure or table would carry it.}
\todo{jakob}{Is this worth a place in the paper? Where?}

\end{document}
